\documentclass[11pt]{article}
\usepackage[a4paper,margin=25mm]{geometry}
\usepackage{amsmath,amssymb,bm,mathtools}

\newcommand{\vect}[1]{\bm{#1}}
\newcommand{\skewm}[1]{[#1]_{\times}}
\newcommand{\pinv}{^{+}}

\begin{document}
\section*{Reusable equations for TRON--AIRBOT cooperative manipulation}

\subsection*{Generalized momentum observer}
\begin{align}
\vect M(\vect q)\dot{\vect v}
+\vect C(\vect q,\vect v)\vect v+\vect g(\vect q)
&=\vect\tau+\vect\tau_{ext}, \tag{OBS-01}\\
\vect\tau_{ext}&=\vect J_e^T\vect h_e, \tag{OBS-02}\\
\vect p&=\vect M(\vect q)\vect v, \tag{OBS-03}\\
\dot{\vect p}&=\vect\tau+\vect\tau_{ext}-\vect\beta,
&\vect\beta&=\vect g-\vect C^T\vect v, \tag{OBS-04}\\
\vect z(t)&=\int_0^t(\vect\tau-\vect\beta+\vect r)\,d\xi, \tag{OBS-05}\\
\vect r&=\vect K_O[(\vect p-\vect p_0)-\vect z], \tag{OBS-06}\\
\dot{\vect r}&=\vect K_O(\vect\tau_{ext}-\vect r). \tag{OBS-07}
\end{align}

\subsection*{Bias compensation and wrench reconstruction}
\begin{align}
\alpha_b&=1-e^{-\Delta t/T_b}, \tag{BIAS-03}\\
\hat{\vect b}_k&=\hat{\vect b}_{k-1}
+\alpha_b(\bar{\vect r}_k-\hat{\vect b}_{k-1}), \tag{BIAS-03b}\\
\hat{\vect\tau}_{ext,k}&=\bar{\vect r}_k-\hat{\vect b}_k, \tag{BIAS-04}\\
\hat{\vect h}_e&=
(\vect J_e\vect J_e^T+\lambda\vect I)^{-1}
\vect J_e\hat{\vect\tau}_{ext}. \tag{WRE-03}
\end{align}

\subsection*{Grasp map and internal wrench}
\begin{align}
\vect G_i&=
\begin{bmatrix}\vect I&\vect0\\\skewm{\vect r_i}&\vect I\end{bmatrix},
&\vect G&=[\vect G_1\ \vect G_2], \tag{GRASP-02}\\
\vect w_o&=\vect G\vect h, \tag{GRASP-03}\\
\vect h_{eff}&=\vect G\pinv\vect w_o,
&\vect h_{int}&=\vect h-\vect h_{eff}, \tag{INT-01}\\
\vect G\vect h_{int}&=\vect0. \tag{INT-02}
\end{align}

\subsection*{Virtual-linkage internal-force coordinates}
\begin{align}
\begin{bmatrix}\vect f_r\\\vect\mu_r\end{bmatrix}
&=\underbrace{\begin{bmatrix}
\vect I&\cdots&\vect I\\
\skewm{\vect r_1}&\cdots&\skewm{\vect r_n}
\end{bmatrix}}_{\vect G_f}\vect f, \tag{VL-02}\\
\vect f&=\vect B_v\vect t+\vect f_e,
&\vect G_f\vect B_v&=\vect0, \tag{VL-03}\\
\hat{\vect t}&=\vect B_v\pinv\vect f, \tag{VL-04}\\
\vect e_{12}&=\frac{\vect p_2-\vect p_1}
{\lVert\vect p_2-\vect p_1\rVert},
&\vect B_{12}&=\begin{bmatrix}\vect e_{12}\\-\vect e_{12}\end{bmatrix}, \tag{VL-05}\\
t_{12}&=\frac12\vect e_{12}^T(\vect f_1-\vect f_2), \tag{VL-06}\\
t_{12}^{*}&=k_t(d_0-d)-d_t\dot d. \tag{VL-07}
\end{align}
For rigid 6D grasps, the complete internal-wrench dimension is
\begin{align}
\dim\ker(\vect G)&=6(n-1), \tag{VL-09}\\
\vect N_G&=\vect I-\vect G\pinv\vect G, \tag{VL-11a}\\
\vect B_{VL}&=\operatorname{orth}(\vect N_G\vect B_{raw}),
&\vect G\vect B_{VL}&=\vect0, \tag{VL-11b}\\
\vect\eta_{VL}&=\vect B_{VL}\pinv\vect h_{int}. \tag{VL-12}
\end{align}

\subsection*{Object impedance and load allocation}
\begin{align}
\vect f_o^*&=\vect K_p(\vect x_o^*-\vect x_o)
+\vect D_p(\dot{\vect x}_o^*-\dot{\vect x}_o)-\hat m_o\vect g, \tag{LOAD-01}\\
\vect h^*&=\vect W^{-1}\vect G^T
(\vect G\vect W^{-1}\vect G^T)\pinv\vect w_o^*. \tag{LOAD-03}
\end{align}

\subsection*{Payload mass and center of mass}
\begin{align}
\vect\theta_o&=[m,mc_x,mc_y,mc_z]^T, \tag{PAY-01}\\
\vect u&=\vect a_o-\vect g, \tag{PAY-02}\\
\vect w_o&=
\underbrace{\begin{bmatrix}
\vect u&\vect0_{3\times3}\\
\vect0_{3\times1}&-\skewm{\vect u}\vect R_{WB}
\end{bmatrix}}_{\vect Y_o}\vect\theta_o. \tag{PAY-04}
\end{align}

\subsection*{CHIP-style hindsight compliance}
\begin{align}
\vect g_{hind}&=\vect g_{ref}-\vect C_f\vect f_{ext}, \tag{CHIP-01}\\
\vect\tau&=\vect g(\vect q)
+\vect J^T[\vect K_x(\vect g-\vect x)-\vect D_x\dot{\vect x}]
+\vect N_J\vect\tau_{posture}, \tag{CHIP-04}\\
\vect K_x&\approx\vect C_f^{-1}, \tag{CHIP-05}\\
\vect N_J&=\vect I-\vect J^T
(\vect J\vect J^T+\lambda_J\vect I)^{-1}\vect J. \tag{CHIP-06}
\end{align}

\clearpage
\section*{Object impedance: derivation and scope}
\noindent Reading note: star2dust, \emph{Impedance control in cooperative
manipulation} (2020-08-17). Equations below are re-derived in the project
convention. They are not a claim of implemented inertia shaping.

\subsection*{Translational target dynamics}
For $\vect e_x=\vect x-\vect x_d$ and positive-definite, constant gains:
\begin{align}
\vect M_d\ddot{\vect e}_x+\vect D_d\dot{\vect e}_x
+\vect K_d\vect e_x&=\vect f_{ext}, \tag{IMP-01}\\
\vect e_{x,ss}&=\vect K_d^{-1}\vect f_{ext},
&\vect C_f&=\vect K_d^{-1}, \tag{IMP-02}\\
\omega_n&=\sqrt{k_d/m_d},
&\zeta&=\frac{d_d}{2\sqrt{m_d k_d}}. \tag{IMP-03}
\end{align}
At a fixed reference, the ideal energy balance is
\begin{equation}
\begin{split}
V&=\tfrac12\dot{\vect e}_x^T\vect M_d\dot{\vect e}_x
+\tfrac12\vect e_x^T\vect K_d\vect e_x,\\
\dot V&=\dot{\vect e}_x^T\vect f_{ext}
-\dot{\vect e}_x^T\vect D_d\dot{\vect e}_x.
\end{split}\tag{IMP-04}
\end{equation}

\subsection*{Object-level wrench command}
World-frame twist is $\vect\nu_o=[\vect v_c^T,\vect\omega^T]^T$;
the reference point is the center of mass.
\begin{equation}
\begin{split}
\vect M_o\dot{\vect\nu}_o+\vect b_o&=\vect w_{ext}+\vect G\vect h,\\
\vect M_o&=\operatorname{diag}(m_o\vect I_3,\vect I_W),\\
\vect b_o&=\begin{bmatrix}-m_o\vect g\\
\vect\omega\times(\vect I_W\vect\omega)\end{bmatrix},\quad
\vect I_W=\vect R_{WB}\vect I_B\vect R_{WB}^T.
\end{split}\tag{IMP-05}
\end{equation}
The following linear pose-error form requires translation or consistent local
small-angle coordinates, not raw Euler-angle derivatives:
\begin{equation}
\vect a_{imp}=\vect a_{ref}+\vect M_d^{-1}
(\vect w_{ext}-\vect D_d\vect e_\nu-\vect K_d\vect e_o).
\tag{IMP-06}
\end{equation}
\begin{equation}
\begin{split}
\vect w_{cmd}={}&\vect b_o+\vect M_o\vect a_{ref}
-\vect M_o\vect M_d^{-1}(\vect D_d\vect e_\nu+\vect K_d\vect e_o)\\
&+(\vect M_o\vect M_d^{-1}-\vect I_6)\vect w_{ext}.
\end{split}\tag{IMP-07}
\end{equation}
\begin{equation}
\vect h_{cmd}=\vect G\pinv\vect w_{cmd}+\vect B_{VL}\vect\eta^*,
\qquad\vect G\vect B_{VL}=\vect0.\tag{IMP-08}
\end{equation}
\noindent Exact wrench realization requires a feasible request and adequate
grasp-map rank. Here $\vect w_{ext}$ excludes robot contact wrenches and
already modeled gravity. See Schneider and Cannon (1992) for the object
impedance framework; the Markdown notes contain source links and caveats.

\clearpage
\section*{Local endpoint impedance: fixed-base derivation}
\noindent $\vect f_i$ denotes force from robot to object, so the reaction
on the robot is $-\vect f_i$. All endpoint equations here are translational.
\begin{equation}
\vect M_{d,i}\ddot{\vect x}_i
+\vect D_{d,i}(\dot{\vect x}_i-\dot{\vect x}_{d,i})
+\vect K_{d,i}(\vect x_i-\vect x_{d,i})=-\vect f_i.
\tag{DIMP-01}
\end{equation}
For a fixed reference and zero initial conditions,
\begin{equation}
\mathcal F_i(s)=(\vect M_{d,i}s^2+\vect D_{d,i}s+\vect K_{d,i})
\mathcal E_i(s),\qquad\vect\epsilon_i=\vect x_{d,i}-\vect x_i.
\tag{DIMP-02}
\end{equation}
\begin{align}
\vect x_{d,i}&=\vect x_i+\vect K_{d,i}^{-1}\vect f_i
\quad\text{(static equilibrium)},\tag{DIMP-03}\\
\vect x_{d,i}^{cmd}&=\vect x_i^{nom}+\vect K_{d,i}^{-1}\vect f_i^*.
\tag{DIMP-04}
\end{align}

\subsection*{Operational-space dynamics}
Assume $\vect J$ has full row rank, and let
$\vect c_q=\vect C_q\dot{\vect q}$.
\begin{equation}
\vect\Lambda=(\vect J\vect M_q^{-1}\vect J^T)^{-1},\qquad
\bar{\vect J}=\vect M_q^{-1}\vect J^T\vect\Lambda.
\tag{DIMP-05}
\end{equation}
\begin{equation}
\begin{split}
\vect\Lambda\ddot{\vect x}+\vect\mu_x+\vect p_x
&=\vect F_{cmd}-\vect f_i,\\
\vect\mu_x&=\vect\Lambda(\vect J\vect M_q^{-1}\vect c_q
-\dot{\vect J}\dot{\vect q}),\\
\vect p_x&=\vect\Lambda\vect J\vect M_q^{-1}\vect g_q,
\quad\vect F_{cmd}=\bar{\vect J}^{T}\vect\tau.
\end{split}\tag{DIMP-06}
\end{equation}
The negative Jacobian-derivative term follows directly from
$\ddot{\vect x}=\vect J\ddot{\vect q}+\dot{\vect J}\dot{\vect q}$.
Define $\vect s_i=\vect D_{d,i}(\dot{\vect x}_i-\dot{\vect x}_{d,i})
+\vect K_{d,i}(\vect x_i-\vect x_{d,i})$. Then
\begin{equation}
\vect F_{cmd,i}=\vect\mu_{x,i}+\vect p_{x,i}
-\vect\Lambda_i\vect M_{d,i}^{-1}\vect s_i
+(\vect I-\vect\Lambda_i\vect M_{d,i}^{-1})\vect f_i.
\tag{DIMP-07}
\end{equation}
\begin{equation}
\vect\tau_i=\vect J_i^T\vect F_{cmd,i}
+(\vect I-\vect J_i^T\bar{\vect J}_i^T)\vect\tau_{0,i}.
\tag{DIMP-08}
\end{equation}
\noindent For an error-acceleration impedance, add reference-acceleration
feedforward. This derivation is for ideal fixed-base torque control without
additional constraint contacts; it is not a floating-base walking controller.
Operational-space background: Khatib (1987).

\clearpage
\section*{Collaborative legged manipulation with PIDC and contact QP}
\noindent Source: Wang et al., \emph{Shared Object Manipulation with a Team
of Collaborative Quadrupeds} (2025). These equations describe the centralized,
torque-controlled formulation. They are planned, not yet implemented for TRON2.

\subsection*{Constrained floating-base dynamics}
For the aggregated generalized coordinates of all robots,
\begin{align}
\vect M\ddot{\vect q}+\vect h
&=\vect S\vect\tau+\vect J_c^T\vect\lambda
+\vect J_x^T\vect F, \tag{PIDC-01}\\
\vect J_c\dot{\vect q}&=\vect0,
&\vect J_c\ddot{\vect q}+\dot{\vect J}_c\dot{\vect q}&=\vect0. \tag{PIDC-02}
\end{align}
The orthogonal projector onto admissible motion and its complementary
constraint projector are
\begin{equation}
\vect P=\vect I-\vect J_c^T(\vect J_c\pinv)^T,\qquad
\vect P^2=\vect P=\vect P^T,\qquad
\vect P\vect J_c^T=\vect0.
\tag{PIDC-03}
\end{equation}
Defining
\begin{equation}
\vect M_c=\vect P\vect M+\vect I-\vect P,\qquad
\bar{\vect M}=\vect M\vect M_c^{-1},
\tag{PIDC-04}
\end{equation}
avoids trying to invert the rank-deficient product $\vect P\vect M$. The
constraint-consistent generalized acceleration is
\begin{equation}
\ddot{\vect q}=\vect M_c^{-1}
\left(\vect P\vect S\vect\tau-\vect P\vect h
+\dot{\vect P}\dot{\vect q}+\vect P\vect J_x^T\vect F\right).
\tag{PIDC-05}
\end{equation}

\subsection*{Constraint-consistent Cartesian impedance}
With $\vect e=\vect x-\vect x_d$, define
\begin{align}
\vect\Lambda_c
&=(\vect J_x\vect M_c^{-1}\vect P\vect J_x^T)^{-1}, \tag{PIDC-06a}\\
\vect h_c
&=\vect\Lambda_c\vect J_x\vect M_c^{-1}
(\vect P\vect h-\dot{\vect P}\dot{\vect q})
-\vect\Lambda_c\dot{\vect J}_x\dot{\vect q}, \tag{PIDC-06b}\\
\vect F_{d,x}
&=\vect h_c+\vect\Lambda_c\ddot{\vect x}_d
-\vect K_d\dot{\vect e}-\vect K_p\vect e, \tag{PIDC-07}\\
\vect\tau_M&=\sum_{x\in\mathcal T}\vect J_x^T\vect F_{d,x},
\quad
\mathcal T=\{\text{object, torsos, swing feet}\}. \tag{PIDC-08}
\end{align}

\subsection*{Grasp closure and object motion}
For contact $i$, $\vect r_i$ points from the object center of mass to the
contact frame and $\vect R_i$ rotates that frame into the object frame:
\begin{align}
\vect G_i&=
\begin{bmatrix}
\vect R_i&\vect0\\
\skewm{\vect r_i}\vect R_i&\vect R_i
\end{bmatrix},
&\vect G&=[\vect G_1\ \cdots\ \vect G_B], \tag{PIDC-09}\\
\vect J_{c,ee}
&=\left[\vect I-\vect G^T(\vect G\pinv)^T\right]\vect J_{ee},
&\vect J_o&=(\vect G\pinv)^T\vect J_{ee}, \tag{PIDC-10}\\
\vect J_c&=
\begin{bmatrix}\vect J_{cf}\\\vect J_{c,ee}\end{bmatrix},
&\vect G\vect\lambda_{ee}^{int}&=\vect0. \tag{PIDC-11}
\end{align}
Thus $\vect J_{c,ee}$ removes relative hand motion that would violate a rigid
grasp, while $\vect J_o$ maps compatible hand motion to object motion.

\subsection*{Underactuation compensation and contact-wrench decision}
The projected inverse-dynamics command is
\begin{align}
\vect B&=\vect I-[(\vect I-\vect S)(\vect I-\vect P)]\pinv,
\tag{PIDC-12a}\\
\vect\tau
&=[\vect P\vect S]\pinv\vect P\vect\tau_M
+\vect B(\vect I-\vect P)\vect\tau_C, \tag{PIDC-12b}\\
(\vect I-\vect P)\vect\tau_C&=\vect J_c^T\vect F_c. \tag{PIDC-12c}
\end{align}
After removing the constant motion term and the orthogonal cross terms from
$\vect\tau^T\vect\tau$, the contact-wrench objective is
\begin{equation}
\min_{\vect F_c}\quad
\vect F_c^T
\underbrace{\vect J_c\vect B^T\vect B\vect J_c^T}_{\vect Q}
\vect F_c.
\tag{PIDC-13}
\end{equation}

\subsection*{Actual contact wrench and physical inequalities}
The optimized command $\vect F_c$ is not generally the actual contact wrench.
Substitution into the constrained dynamics gives
\begin{equation}
\vect\lambda=\vect\rho\vect F_c+\vect\eta,
\tag{PIDC-14}
\end{equation}
where
\begin{equation}
\begin{split}
\vect\rho={}&(\vect J_c^T)\pinv(\vect I-\vect P)
(\bar{\vect M}\vect P-\vect I)\vect S\vect B\vect J_c^T,\\
\vect\eta={}&(\vect J_c^T)\pinv(\vect I-\vect P)
\bigl[\vect\varepsilon\vect\tau_M
+\bar{\vect M}(-\vect P\vect h+\dot{\vect P}\dot{\vect q})+\vect h\\
&\hspace{37mm}
+(\bar{\vect M}\vect P-\vect I)\vect J_x^T\vect F\bigr],\\
\vect\varepsilon={}&(\bar{\vect M}\vect P-\vect I)
\vect S[\vect P\vect S]\pinv\vect P.
\end{split}
\tag{PIDC-15}
\end{equation}
For actual point-contact force $\vect\lambda_{f,i}$, with tangent directions
$\vect n_{x,i},\vect n_{y,i}$, normal $\vect n_{z,i}$, and friction $\mu_i$,
\begin{equation}
\underbrace{\begin{bmatrix}-\infty\\-\infty\\0\\0\\0\end{bmatrix}}_{\underline{\vect d}}
\le
\underbrace{\begin{bmatrix}
(\vect n_{x,i}-\mu_i\vect n_{z,i})^T\\
(\vect n_{y,i}-\mu_i\vect n_{z,i})^T\\
(\vect n_{x,i}+\mu_i\vect n_{z,i})^T\\
(\vect n_{y,i}+\mu_i\vect n_{z,i})^T\\
\vect n_{z,i}^T
\end{bmatrix}}_{\vect C_i}
\vect\lambda_{f,i}
\le
\underbrace{\begin{bmatrix}0\\0\\+\infty\\+\infty\\+\infty\end{bmatrix}}_{\bar{\vect d}}.
\tag{PIDC-16}
\end{equation}
For a rectangular surface contact
$\vect\lambda_i=[f_x,f_y,f_z,m_x,m_y,m_z]^T$, rolling and torsional limits are
\begin{equation}
\vect0\le
\underbrace{\begin{bmatrix}
0&0&\xi&0&0&-1\\
0&0&\xi&0&0& 1\\
0&0&\delta_x&-1&0&0\\
0&0&\delta_x& 1&0&0\\
0&0&\delta_y&0&-1&0\\
0&0&\delta_y&0& 1&0
\end{bmatrix}}_{\vect C_{m,i}}
\vect\lambda_i.
\tag{PIDC-17}
\end{equation}
Equivalently,
$|m_z|\le\xi f_z$, $|m_x|\le\delta_xf_z$, and
$|m_y|\le\delta_yf_z$.

Let $\vect S_{ns}$ select only actuated joints. The actuator bounds become
linear inequalities in $\vect F_c$:
\begin{equation}
\underline{\vect\tau}\le\vect\Xi\vect F_c\le\bar{\vect\tau},
\quad
\vect\Xi=\vect S_{ns}^T\vect B\vect J_c^T,
\tag{PIDC-18}
\end{equation}
\begin{equation}
\begin{split}
\underline{\vect\tau}&=\vect\tau_{min}
-\vect S_{ns}^T[\vect P\vect S]\pinv\vect P\vect\tau_M,\\
\bar{\vect\tau}&=\vect\tau_{max}
-\vect S_{ns}^T[\vect P\vect S]\pinv\vect P\vect\tau_M.
\end{split}
\tag{PIDC-19}
\end{equation}

\subsection*{Standard QP and final command}
For selection matrices $\vect E_i$ extracting contact $i$ from
$\vect\lambda$, substitute
$\vect\lambda_i=\vect E_i(\vect\rho\vect F_c+\vect\eta)$ into
the friction and moment inequalities. Stack them with the actuator bounds:
\begin{equation}
\begin{aligned}
\min_{\vect F_c}\quad&
\tfrac12\vect F_c^T\vect H\vect F_c+\vect g^T\vect F_c,\\
\text{s.t.}\quad&\vect l\le\vect A\vect F_c\le\vect u,\\
\vect H&=2(\vect J_c\vect B^T\vect B\vect J_c^T
+\epsilon_Q\vect I),\qquad \vect g=\vect0.
\end{aligned}
\tag{PIDC-20}
\end{equation}
The small $\epsilon_Q>0$ is an implementation regularizer, not part of the
paper's original objective. After solving the QP,
\begin{align}
\vect\tau^*&=[\vect P\vect S]\pinv\vect P\vect\tau_M
+\vect B\vect J_c^T\vect F_c^*, \tag{PIDC-21a}\\
\vect\lambda^*&=\vect\rho\vect F_c^*+\vect\eta. \tag{PIDC-21b}
\end{align}

\subsection*{TRON2 scope warning}
The paper assumes stationary point feet, $\vect J_{cf}\dot{\vect q}=0$.
For TRON2 wheelfoot locomotion, rolling velocity is allowed along the wheel
heading; $\vect J_{cf}$ must therefore be replaced by nonholonomic rolling
constraints. The stationary form applies directly only to solefoot or parked
support. A full implementation also requires a verified joint-torque interface
and a grasp model consistent with the wrench directions actually supported by
the AIRBOT gripper.

\end{document}
